% =========================================================
% Proof Stub: Even-Odd Convergence
% Source: volume-iii/analysis/sequences/notes/sequences/notes-subsequence-toolkit.tex
% Mode: standard
% =========================================================

\clearpage
\phantomsection
\hypertarget{proof-RA-ABB-C-even-odd-conv}{}

\subsubsection[Even-Odd Convergence]{Proof --- Even-Odd Convergence}

\bigskip

\noindent
\textbf{Source.}
See \texttt{notes-subsequence-toolkit.tex}.

\vspace{0.75em}

\noindent
\textbf{Goal.}
If $(a_{2n}) \to L$ and $(a_{2n+1}) \to L$, then $(a_n) \to L$.

\vspace{0.75em}

\noindent
\textbf{Logical form.}
\[
  (a_{2n} \to L) \wedge (a_{2n+1} \to L) \Rightarrow a_n \to L.
\]

\vspace{0.75em}

\noindent
\textbf{Proof strategy.}
Residue-Class Convergence with $k = 2$.

\vspace{0.75em}

\noindent
\textbf{Proof.}
\begin{proof}
\Claim If $(a_{2n}) \to L$ and $(a_{2n+1}) \to L$, then $(a_n) \to L$.

\Given 

\Goal 

\Strategy Residue-Class Convergence with $k = 2$.

\medskip

% --- apply Residue-Class Convergence with k = 2 ---

\end{proof}

\begin{remark*}[Proof shape]
Special case $k=2$ of Residue-Class Convergence.
\end{remark*}

\begin{remark*}[Dependencies]
\begin{itemize}
  \item Residue-Class Convergence.
\end{itemize}
\end{remark*}
